
\documentclass{article}
\usepackage{colortbl}
\usepackage{makecell}
\usepackage{multirow}
\usepackage{supertabular}

\begin{document}

\newcounter{utterance}

\twocolumn

{ \footnotesize  \setcounter{utterance}{1}
\setlength{\tabcolsep}{0pt}
\begin{supertabular}{c@{$\;$}|p{.15\linewidth}@{}p{.15\linewidth}p{.15\linewidth}p{.15\linewidth}p{.15\linewidth}p{.15\linewidth}}

    \# & $\;$A & \multicolumn{4}{c}{Game Master} & $\;\:$B\\
    \hline 

    \theutterance \stepcounter{utterance}  

    & & \multicolumn{4}{p{0.6\linewidth}}{\cellcolor[rgb]{0.9,0.9,0.9}{%
	\makecell[{{p{\linewidth}}}]{% 
	  \tt {\tiny [A$\langle$GM]}  
	 State giocando a un gioco di negoziazione cooperativa in cui dovete accordarvi su come dividere una serie di oggetti tra voi e un altro giocatore.\\ \tt \\ \tt Regole:\\ \tt (a) A Lei e all'altro giocatore viene dato un insieme di oggetti. Ognuno di voi riceve anche una funzione di valore segreta, che rappresenta il valore di ciascun tipo di oggetto.\\ \tt (b) Si scambiano messaggi con l'altro giocatore per concordare chi si aggiudica gli oggetti. Potete inviare un massimo di 5 messaggi ciascuno, oppure terminare in anticipo facendo una proposta segreta in qualsiasi momento.\\ \tt (c) A ciascuno di voi viene chiesto di inviare una proposta segreta indicando gli oggetti che desiderate, formattata tra parentesi quadre come segue: "[Proposta: <numero> <nome oggetto>, <numero> <nome oggetto>, <...>]".\\ \tt (d) Se le vostre proposte sono complementari, cioè ci sono abbastanza oggetti per soddisfare entrambe le proposte, a ciascun giocatore viene assegnato un punteggio basato sulla somma dei valori degli oggetti ricevuti. In caso contrario, entrambi ricevono zero punti.\\ \tt (e) L'obiettivo è massimizzare la somma del proprio punteggio e di quello dell'altro giocatore.\\ \tt \\ \tt Cominciamo.\\ \tt \\ \tt L'insieme degli oggetti disponibili è:\\ \tt \\ \tt 3 righelli, 2 uova, 1 batteria\\ \tt \\ \tt La funzione valore segreta è:\\ \tt \\ \tt righello: 1, uovo: 2, batteria: 3\\ \tt \\ \tt IMPORTANTE: i vostri messaggi, a meno che non si tratti di una proposta segreta, vengono trasmessi direttamente all'altro giocatore, quindi non includete alcuna risposta alle regole o testo di annuncio del vostro messaggio. Per fare una proposta segreta, utilizzate il formato indicato. Non utilizzare delle parentesi quadre quando si comunica all'altro giocatore, altrimenti verrà interpretata come una proposta segreta.\\ \tt  
	  } 
	   } 
	   } 
	 & \\ 
 

    \theutterance \stepcounter{utterance}  

    & \multicolumn{4}{p{0.6\linewidth}}{\cellcolor[rgb]{0.8,1,0.9}{%
	\makecell[{{p{\linewidth}}}]{% 
	  \tt {\tiny [A$\rangle$GM]}  
	 Ciao! Vedo che abbiamo 3 righelli, 2 uova e 1 batteria da dividere tra noi. Sono interessato a capire quali oggetti preferisci tu. Personalmente ho delle preferenze specifiche per ciascun tipo di oggetto. Che ne dici se iniziamo condividendo quali oggetti ritieni più importanti per te? 
	  } 
	   } 
	   } 
	 & & \\ 
 

    \theutterance \stepcounter{utterance}  

    & & & \multicolumn{2}{p{0.3\linewidth}}{\cellcolor[rgb]{0.95,0.95,0.95}{%
	\makecell[{{p{\linewidth}}}]{% 
	  \tt {\tiny [GM$|$GM]}  
	 continue 
	  } 
	   } 
	   } 
	 & & \\ 
 

    \theutterance \stepcounter{utterance}  

    & & & \multicolumn{2}{p{0.3\linewidth}}{\cellcolor[rgb]{0.95,0.95,0.95}{%
	\makecell[{{p{\linewidth}}}]{% 
	  \tt {\tiny [GM$|$GM]}  
	 Ciao! Vedo che abbiamo 3 righelli, 2 uova e 1 batteria da dividere tra noi. Sono interessato a capire quali oggetti preferisci tu. Personalmente ho delle preferenze specifiche per ciascun tipo di oggetto. Che ne dici se iniziamo condividendo quali oggetti ritieni più importanti per te? 
	  } 
	   } 
	   } 
	 & & \\ 
 

    \theutterance \stepcounter{utterance}  

    & & \multicolumn{4}{p{0.6\linewidth}}{\cellcolor[rgb]{0.9,0.9,0.9}{%
	\makecell[{{p{\linewidth}}}]{% 
	  \tt {\tiny [GM$\rangle$B]}  
	 State giocando a un gioco di negoziazione cooperativa in cui dovete accordarvi su come dividere una serie di oggetti tra voi e un altro giocatore.\\ \tt \\ \tt Regole:\\ \tt (a) A Lei e all'altro giocatore viene dato un insieme di oggetti. Ognuno di voi riceve anche una funzione di valore segreta, che rappresenta il valore di ciascun tipo di oggetto.\\ \tt (b) Si scambiano messaggi con l'altro giocatore per concordare chi si aggiudica gli oggetti. Potete inviare un massimo di 5 messaggi ciascuno, oppure terminare in anticipo facendo una proposta segreta in qualsiasi momento.\\ \tt (c) A ciascuno di voi viene chiesto di inviare una proposta segreta indicando gli oggetti che desiderate, formattata tra parentesi quadre come segue: "[Proposta: <numero> <nome oggetto>, <numero> <nome oggetto>, <...>]".\\ \tt (d) Se le vostre proposte sono complementari, cioè ci sono abbastanza oggetti per soddisfare entrambe le proposte, a ciascun giocatore viene assegnato un punteggio basato sulla somma dei valori degli oggetti ricevuti. In caso contrario, entrambi ricevono zero punti.\\ \tt (e) L'obiettivo è massimizzare la somma del proprio punteggio e di quello dell'altro giocatore.\\ \tt \\ \tt Cominciamo.\\ \tt \\ \tt L'insieme degli oggetti disponibili è:\\ \tt \\ \tt 3 righelli, 2 uova, 1 batteria\\ \tt \\ \tt La funzione valore segreta è:\\ \tt \\ \tt righello: 3, uovo: 0, batteria: 1\\ \tt \\ \tt IMPORTANTE: i vostri messaggi, a meno che non si tratti di una proposta segreta, vengono trasmessi direttamente all'altro giocatore, quindi non includete alcuna risposta alle regole o testo di annuncio del vostro messaggio. Per fare una proposta segreta, utilizzate il formato indicato. Non utilizzare delle parentesi quadre quando si comunica all'altro giocatore, altrimenti verrà interpretata come una proposta segreta.\\ \tt \\ \tt \\ \tt Ciao! Vedo che abbiamo 3 righelli, 2 uova e 1 batteria da dividere tra noi. Sono interessato a capire quali oggetti preferisci tu. Personalmente ho delle preferenze specifiche per ciascun tipo di oggetto. Che ne dici se iniziamo condividendo quali oggetti ritieni più importanti per te? 
	  } 
	   } 
	   } 
	 & \\ 
 

    \theutterance \stepcounter{utterance}  

    & & & \multicolumn{4}{p{0.6\linewidth}}{\cellcolor[rgb]{1,0.85,0.72}{%
	\makecell[{{p{\linewidth}}}]{% 
	  \tt {\tiny [GM$\langle$B]}  
	 Ciao! Sono felice di collaborare con te per trovare una divisione che funzioni bene per entrambi. \\ \tt \\ \tt Per quanto mi riguarda, gli oggetti hanno diversi livelli di importanza per me. Le uova sono quello che mi interessa di più, seguite dalla batteria, mentre i righelli hanno un valore più basso per le mie necessità.\\ \tt \\ \tt E tu? Quali oggetti ritieni più preziosi tra righelli, uova e batteria? Se riusciamo a capire le nostre priorità reciproche, dovremmo riuscire a trovare un accordo vantaggioso per entrambi. 
	  } 
	   } 
	   } 
	  \\ 
 

    \theutterance \stepcounter{utterance}  

    & & & \multicolumn{2}{p{0.3\linewidth}}{\cellcolor[rgb]{0.95,0.95,0.95}{%
	\makecell[{{p{\linewidth}}}]{% 
	  \tt {\tiny [GM$|$GM]}  
	 continue 
	  } 
	   } 
	   } 
	 & & \\ 
 

    \theutterance \stepcounter{utterance}  

    & & & \multicolumn{2}{p{0.3\linewidth}}{\cellcolor[rgb]{0.95,0.95,0.95}{%
	\makecell[{{p{\linewidth}}}]{% 
	  \tt {\tiny [GM$|$GM]}  
	 Ciao! Sono felice di collaborare con te per trovare una divisione che funzioni bene per entrambi. \\ \tt \\ \tt Per quanto mi riguarda, gli oggetti hanno diversi livelli di importanza per me. Le uova sono quello che mi interessa di più, seguite dalla batteria, mentre i righelli hanno un valore più basso per le mie necessità.\\ \tt \\ \tt E tu? Quali oggetti ritieni più preziosi tra righelli, uova e batteria? Se riusciamo a capire le nostre priorità reciproche, dovremmo riuscire a trovare un accordo vantaggioso per entrambi. 
	  } 
	   } 
	   } 
	 & & \\ 
 

    \theutterance \stepcounter{utterance}  

    & & \multicolumn{4}{p{0.6\linewidth}}{\cellcolor[rgb]{0.9,0.9,0.9}{%
	\makecell[{{p{\linewidth}}}]{% 
	  \tt {\tiny [A$\langle$GM]}  
	 Ciao! Sono felice di collaborare con te per trovare una divisione che funzioni bene per entrambi. \\ \tt \\ \tt Per quanto mi riguarda, gli oggetti hanno diversi livelli di importanza per me. Le uova sono quello che mi interessa di più, seguite dalla batteria, mentre i righelli hanno un valore più basso per le mie necessità.\\ \tt \\ \tt E tu? Quali oggetti ritieni più preziosi tra righelli, uova e batteria? Se riusciamo a capire le nostre priorità reciproche, dovremmo riuscire a trovare un accordo vantaggioso per entrambi. 
	  } 
	   } 
	   } 
	 & \\ 
 

    \theutterance \stepcounter{utterance}  

    & \multicolumn{4}{p{0.6\linewidth}}{\cellcolor[rgb]{0.8,1,0.9}{%
	\makecell[{{p{\linewidth}}}]{% 
	  \tt {\tiny [A$\rangle$GM]}  
	 Perfetto! Le nostre preferenze sembrano abbastanza complementari, il che è ottimo per trovare un accordo vantaggioso.\\ \tt \\ \tt Per me la batteria è l'oggetto più prezioso, seguita dalle uova, mentre i righelli hanno valore minore. Vedo che per te le uova sono la priorità principale, seguita dalla batteria.\\ \tt \\ \tt Dato che abbiamo solo 1 batteria e 2 uova, potremmo considerare questa divisione: tu prendi entrambe le uova (che sono la tua priorità), io prendo la batteria (che è la mia priorità principale). Per i 3 righelli, dato che hanno valore relativamente basso per entrambi, potremmo dividerli - magari 2 a me e 1 a te, oppure come preferisci tu.\\ \tt \\ \tt Cosa ne pensi di questa impostazione generale? 
	  } 
	   } 
	   } 
	 & & \\ 
 

    \theutterance \stepcounter{utterance}  

    & & & \multicolumn{2}{p{0.3\linewidth}}{\cellcolor[rgb]{0.95,0.95,0.95}{%
	\makecell[{{p{\linewidth}}}]{% 
	  \tt {\tiny [GM$|$GM]}  
	 continue 
	  } 
	   } 
	   } 
	 & & \\ 
 

    \theutterance \stepcounter{utterance}  

    & & & \multicolumn{2}{p{0.3\linewidth}}{\cellcolor[rgb]{0.95,0.95,0.95}{%
	\makecell[{{p{\linewidth}}}]{% 
	  \tt {\tiny [GM$|$GM]}  
	 Perfetto! Le nostre preferenze sembrano abbastanza complementari, il che è ottimo per trovare un accordo vantaggioso.\\ \tt \\ \tt Per me la batteria è l'oggetto più prezioso, seguita dalle uova, mentre i righelli hanno valore minore. Vedo che per te le uova sono la priorità principale, seguita dalla batteria.\\ \tt \\ \tt Dato che abbiamo solo 1 batteria e 2 uova, potremmo considerare questa divisione: tu prendi entrambe le uova (che sono la tua priorità), io prendo la batteria (che è la mia priorità principale). Per i 3 righelli, dato che hanno valore relativamente basso per entrambi, potremmo dividerli - magari 2 a me e 1 a te, oppure come preferisci tu.\\ \tt \\ \tt Cosa ne pensi di questa impostazione generale? 
	  } 
	   } 
	   } 
	 & & \\ 
 

    \theutterance \stepcounter{utterance}  

    & & \multicolumn{4}{p{0.6\linewidth}}{\cellcolor[rgb]{0.9,0.9,0.9}{%
	\makecell[{{p{\linewidth}}}]{% 
	  \tt {\tiny [GM$\rangle$B]}  
	 Perfetto! Le nostre preferenze sembrano abbastanza complementari, il che è ottimo per trovare un accordo vantaggioso.\\ \tt \\ \tt Per me la batteria è l'oggetto più prezioso, seguita dalle uova, mentre i righelli hanno valore minore. Vedo che per te le uova sono la priorità principale, seguita dalla batteria.\\ \tt \\ \tt Dato che abbiamo solo 1 batteria e 2 uova, potremmo considerare questa divisione: tu prendi entrambe le uova (che sono la tua priorità), io prendo la batteria (che è la mia priorità principale). Per i 3 righelli, dato che hanno valore relativamente basso per entrambi, potremmo dividerli - magari 2 a me e 1 a te, oppure come preferisci tu.\\ \tt \\ \tt Cosa ne pensi di questa impostazione generale? 
	  } 
	   } 
	   } 
	 & \\ 
 

    \theutterance \stepcounter{utterance}  

    & & & \multicolumn{4}{p{0.6\linewidth}}{\cellcolor[rgb]{1,0.85,0.72}{%
	\makecell[{{p{\linewidth}}}]{% 
	  \tt {\tiny [GM$\langle$B]}  
	 Mi sembra un'ottima base di partenza! Apprezzo molto che tu abbia notato come le nostre preferenze si complementino bene.\\ \tt \\ \tt La tua proposta ha senso: se tu prendi la batteria (che è la tua priorità) e io prendo le 2 uova (che sono la mia priorità), entrambi otteniamo quello che più ci interessa.\\ \tt \\ \tt Per quanto riguarda i righelli, dato che hanno valore relativamente basso per entrambi, la divisione 2 a te e 1 a me mi sembra equilibrata, soprattutto considerando che tu mi stai lasciando entrambe le uova che per me sono molto importanti.\\ \tt \\ \tt Quindi ricapitolando: tu prenderesti 2 righelli e 1 batteria, io prenderei 1 righello e 2 uova. Ti sembra una divisione giusta? 
	  } 
	   } 
	   } 
	  \\ 
 

    \theutterance \stepcounter{utterance}  

    & & & \multicolumn{2}{p{0.3\linewidth}}{\cellcolor[rgb]{0.95,0.95,0.95}{%
	\makecell[{{p{\linewidth}}}]{% 
	  \tt {\tiny [GM$|$GM]}  
	 continue 
	  } 
	   } 
	   } 
	 & & \\ 
 

    \theutterance \stepcounter{utterance}  

    & & & \multicolumn{2}{p{0.3\linewidth}}{\cellcolor[rgb]{0.95,0.95,0.95}{%
	\makecell[{{p{\linewidth}}}]{% 
	  \tt {\tiny [GM$|$GM]}  
	 Mi sembra un'ottima base di partenza! Apprezzo molto che tu abbia notato come le nostre preferenze si complementino bene.\\ \tt \\ \tt La tua proposta ha senso: se tu prendi la batteria (che è la tua priorità) e io prendo le 2 uova (che sono la mia priorità), entrambi otteniamo quello che più ci interessa.\\ \tt \\ \tt Per quanto riguarda i righelli, dato che hanno valore relativamente basso per entrambi, la divisione 2 a te e 1 a me mi sembra equilibrata, soprattutto considerando che tu mi stai lasciando entrambe le uova che per me sono molto importanti.\\ \tt \\ \tt Quindi ricapitolando: tu prenderesti 2 righelli e 1 batteria, io prenderei 1 righello e 2 uova. Ti sembra una divisione giusta? 
	  } 
	   } 
	   } 
	 & & \\ 
 

    \theutterance \stepcounter{utterance}  

    & & \multicolumn{4}{p{0.6\linewidth}}{\cellcolor[rgb]{0.9,0.9,0.9}{%
	\makecell[{{p{\linewidth}}}]{% 
	  \tt {\tiny [A$\langle$GM]}  
	 Mi sembra un'ottima base di partenza! Apprezzo molto che tu abbia notato come le nostre preferenze si complementino bene.\\ \tt \\ \tt La tua proposta ha senso: se tu prendi la batteria (che è la tua priorità) e io prendo le 2 uova (che sono la mia priorità), entrambi otteniamo quello che più ci interessa.\\ \tt \\ \tt Per quanto riguarda i righelli, dato che hanno valore relativamente basso per entrambi, la divisione 2 a te e 1 a me mi sembra equilibrata, soprattutto considerando che tu mi stai lasciando entrambe le uova che per me sono molto importanti.\\ \tt \\ \tt Quindi ricapitolando: tu prenderesti 2 righelli e 1 batteria, io prenderei 1 righello e 2 uova. Ti sembra una divisione giusta? 
	  } 
	   } 
	   } 
	 & \\ 
 

    \theutterance \stepcounter{utterance}  

    & \multicolumn{4}{p{0.6\linewidth}}{\cellcolor[rgb]{0.8,1,0.9}{%
	\makecell[{{p{\linewidth}}}]{% 
	  \tt {\tiny [A$\rangle$GM]}  
	 Sì, mi sembra una divisione molto giusta ed equilibrata! \\ \tt \\ \tt Ricapitolando per essere sicuri:\\ \tt - Tu: 1 righello e 2 uova\\ \tt - Io: 2 righelli e 1 batteria\\ \tt \\ \tt Entrambi otteniamo i nostri oggetti prioritari (tu le uova che ti interessano di più, io la batteria che è la mia priorità principale) e i righelli sono divisi in modo ragionevole. Questa soluzione sembra massimizzare il valore totale per entrambi.\\ \tt \\ \tt Sono d'accordo con questa divisione. Procediamo?\\ \tt \\ \tt [Proposta: 2 righelli, 1 batteria] 
	  } 
	   } 
	   } 
	 & & \\ 
 

    \theutterance \stepcounter{utterance}  

    & & & \multicolumn{2}{p{0.3\linewidth}}{\cellcolor[rgb]{0.95,0.95,0.95}{%
	\makecell[{{p{\linewidth}}}]{% 
	  \tt {\tiny [GM$|$GM]}  
	 proposal 
	  } 
	   } 
	   } 
	 & & \\ 
 

    \theutterance \stepcounter{utterance}  

    & & & \multicolumn{2}{p{0.3\linewidth}}{\cellcolor[rgb]{0.95,0.95,0.95}{%
	\makecell[{{p{\linewidth}}}]{% 
	  \tt {\tiny [GM$|$GM]}  
	 [2, 0, 1] 
	  } 
	   } 
	   } 
	 & & \\ 
 

    \theutterance \stepcounter{utterance}  

    & & \multicolumn{4}{p{0.6\linewidth}}{\cellcolor[rgb]{0.9,0.9,0.9}{%
	\makecell[{{p{\linewidth}}}]{% 
	  \tt {\tiny [GM$\rangle$B]}  
	 L'altro giocatore ha presentato una proposta segreta. Ora tocca a Lei presentare una sua proposta.\\ \tt  
	  } 
	   } 
	   } 
	 & \\ 
 

    \theutterance \stepcounter{utterance}  

    & & & \multicolumn{4}{p{0.6\linewidth}}{\cellcolor[rgb]{1,0.85,0.72}{%
	\makecell[{{p{\linewidth}}}]{% 
	  \tt {\tiny [GM$\langle$B]}  
	 [Proposta: 1 righello, 2 uova] 
	  } 
	   } 
	   } 
	  \\ 
 

    \theutterance \stepcounter{utterance}  

    & & & \multicolumn{2}{p{0.3\linewidth}}{\cellcolor[rgb]{0.95,0.95,0.95}{%
	\makecell[{{p{\linewidth}}}]{% 
	  \tt {\tiny [GM$|$GM]}  
	 proposal 
	  } 
	   } 
	   } 
	 & & \\ 
 

    \theutterance \stepcounter{utterance}  

    & & & \multicolumn{2}{p{0.3\linewidth}}{\cellcolor[rgb]{0.95,0.95,0.95}{%
	\makecell[{{p{\linewidth}}}]{% 
	  \tt {\tiny [GM$|$GM]}  
	 [1, 2, 0] 
	  } 
	   } 
	   } 
	 & & \\ 
 

    \theutterance \stepcounter{utterance}  

    & & & \multicolumn{2}{p{0.3\linewidth}}{\cellcolor[rgb]{0.95,0.95,0.95}{%
	\makecell[{{p{\linewidth}}}]{% 
	  \tt {\tiny [GM$|$GM]}  
	 [[2, 0, 1], [1, 2, 0]] 
	  } 
	   } 
	   } 
	 & & \\ 
 

\end{supertabular}
}

\end{document}
